\documentclass[10pt, a4paper]{article}
\usepackage{preamble}

\title{Algebra II \\
    \large Problem Class}
\author{Luke Phillips}
\date{October 2025}

\begin{document}

\maketitle

\newpage

\tableofcontents

\newpage

\section{Problems Class \texorpdfstring{$1$}{}}

\begin{problem}[Question $2$]
    Let $R$ be a ring.
    Show,
    using only the definition of ring,
    that for any two elements $a, b \in R$ we have
    \begin{enumerate}[label = (\alph*)]
        \item $0 \cdot a = a \cdot 0 = 0$.

        \item $(-a)b = -(ab)$,
        and similarly $a(-b) = -(ab)$.

        \item $(-a)(-b) = ab$
    \end{enumerate}

    \begin{solution}
        \begin{enumerate}[label = (\alph*)]
            \item
            \begin{align*}
                0 \cdot a &= (0 + 0) \cdot a &&\text{(by adding $0$}) \\
                &= 0 \cdot a + 0 \cdot a &&\text{(by distributivity)} \\
                &\iff \\
                0 &= 0 \cdot a &&\text{(subtracting $0 \cdot a$)}
            \end{align*}
            $a \cdot 0$ is done similarly.

            \item
            First we want:
            \begin{align*}
                (-a)b &= -(ab) \\
                &\iff \\
                (-a)b + ab &= 0 &&\text{(by definition of $-x$)} \\
                &\iff \\
                (-a + a)b &= 0  &&\text{(by distributivity)} \\
                &\iff \\
                0 \cdot b &= 0
            \end{align*}
            this is true by $a$.

            \item
            Sub $-b$ in $b$:
            \[
            (-a)(-b) = -(a(-b))
            \]
            we want $-(a(-b)) = ab$,
            i.e.
            $ab + a(-b) = 0$.
            But $ab + a(-b) = a(b + -b) = a \cdot 0 = 0$ by $a$.
        \end{enumerate}
    \end{solution}
\end{problem}

\begin{problem}[Question $3$]
    Let $V$ be a finite dimensional vector space over $\C$.
    Let
    \[
    \mathrm{End}(V) = \{f : V \to V \mid f\text{ is a linear map}\}.
    \]
    An element of $\mathrm{End}(V)$ is usually called a
    (linear)
    endomorphism.
    We can add two such functions:
    $(f + g)(v) = f(v) + g(v)$,
    for $f, g \in \mathrm{End}(V)$,
    $v \in V$.
    Define another binary operation by composition:
    \[
    (f \circ g)(v) = f(g(v)).
    \]
    Prove that $(\mathrm{End}(V), +, \circ)$ is a ring.

    \begin{solution}
        First we know that the endomorphism form an abelian group under addition.
        Identity:
        $\mathrm{id} : V \to V$,
        $V \mapsto V$.
        Associativity:
        We need,
        $\forall f, g, h \in \mathrm{End}(V)$,
        \[
        f \circ (g \circ h) = (f \circ g) \circ h.
        \]
        For any $v \in V$,
        \[
        (f \circ (g \circ h))(v) = f((g \circ h)(v)) = f(g(h(v)))
        \]
        and $((f \circ g) \circ h)(v) = (f \circ g)(h(v)) = f(g(h(v)))$.
        This proves associativity.
        This works for any functions $f, g, h$ $X \to X$,
        $X$ any set.

        Distributivity:
        \begin{align*}
            (f \circ (g + h))(v) &= f(g(v) + h(v)) \\
            &= f(g(v)) + f(h(v)) &&\text{(by linearity of $f$)} \\
            &= (f \circ g)(v) + (f \circ h)(v)
        \end{align*}
        Similarly for right distributivity.
        Thus,
        $(\mathrm{End}(V), +, \circ)$ is a ring.
    \end{solution}
\end{problem}

\begin{problem}[Question $17$]
    \begin{enumerate}[label = (\alph*)]
        \item Let $m, n, a \in \Z$ be such that $m$ and $n$ are coprime.
        Prove that if $m \mid na$,
        then $m \mid a$.

        \item Show,
        by giving a counter-example,
        that the statement in the previous part fails if $m$ and $n$ are not coprime.

        \item Let $m, n \in \Z$ be coprime.
        Prove that if $m \mid a$ and $n \mid a$,
        for some integer $a$,
        then
        \[
        mn \mid a.
        \]
    \end{enumerate}

    \begin{solution}
        \begin{enumerate}[label = (\alph*)]
            \item
            Coprime when $\gcd(m, n) = 1$.
            By the Euclidean algorithm,
            $\exists x, y \in \Z$ such that $xm + yn = 1$.
            $m \mid na \iff na = mr$,
            some $r \in \Z$.
            By the Euclidean algorithm equation,
            we get
            \[
            axm + ayn = a \iff axm + mry = a \iff m(ax + ry) = a \implies m \mid a.
            \]

            \item $m = 2$,
            $n = 4$.

            \item
            Write $a = mr = ns$,
            for some $r, s \in \Z$.
            Then $m \mid ns$.
            But by ($a$)
            $m \mid s$,
            i.e.,
            $s = mt$,
            for some $t \in \Z$.
            Thus $a = ns = nmt$,
            i.e. $nm \mid a$.
        \end{enumerate}
    \end{solution}
\end{problem}

\begin{problem}[Question $21$]
    Let
    \[
    f(x) = x ^ 3 - 12x ^ 2 - 42,\qquad g(x) = x - 3
    \]
    be polynomials in $\Q[x]$.
    Find $q(x)$ and $r(x)$ in $\Q[x]$ such that $\deg{r} < 1$ and
    \[
    f(x) = q(x)g(x) + r(x).
    \]

    \begin{solution}
        \[
        -9x ^ 2 + 27x +39.
        \]
    \end{solution}
\end{problem}

\newpage

\section{Problems Class \texorpdfstring{$2$}{}}

\begin{problem}[Question $6$]
    \begin{enumerate}[label = (\alph*)]
        \item Show elements $3$,
        $1 - \sqrt{-5}$ and $1 + \sqrt{-5}$ are irreducible in $\Z[\sqrt{-5}]$.

        \item Show that $1 - \sqrt{-5}$ is not a prime element in $\Z[-\sqrt{5}]$.
    \end{enumerate}

    \begin{solution}
        \begin{enumerate}[label = (\alph*)]
            \item
            Norm map:
            $N : \Z[\sqrt{-5}] \to \N \cup \{0\}$,
            $N(a + b\sqrt{-5}) + a ^ 2 + 5b ^ 2 = (a + b\sqrt{-5})(a - b\sqrt{5}) = 2\bar{z}$.

            Fact:
            $N$ is multiplicative,
            i.e.,
            $N(xy) = N(x)N(y)$,
            since
            \[
            N(xy) = xy\cdot \overline{xy} = xy \cdot \bar{x} \cdot \bar{y} = x\bar{x} \cdot y\bar{y} = N(x)N(y).
            \]

            Assume $3$ has a factorisation:
            \[
            3 = (x + y\sqrt{-5})(z + w\sqrt{-5}).
            \]
            Apply $N$:
            \[
            N(3) = 9 = (x ^ 2 + 5y ^ 2)(z ^ 2 + 5w ^ 2)
            \]
            This means that $x ^ 2 + 5y ^ 2 \mid 9$,
            so $x ^ 2 + 5y ^ 2 \in \{\pm 1, \pm 3, \pm 9\}$ cannot have $-1, -3, -9$.
            if $x ^ 2 + 5y ^ 2 = 1$ then $y = 0$
            (if $y \geq 1$,
            then $x ^ 2 + 5y ^ 2 \geq 5$)
            thus $x = \pm 1$.
            If $x ^ 2 + 5y ^ 2 = 3$ then $y = 0$ again $x ^ 2 = 3$ has no solution.
            If $x ^ 2 + 5y ^ 2 = 9$ then $y = 0$,
            $x = \pm 3$ or $y = \pm 1$,
            $x = \pm 2$
            \[
            x + y\sqrt{-5} = \pm 3 \implies z + w\sqrt{-5} = \pm 1
            \]
            a unit.
            \[
            x + y\sqrt{-5} = \pm 2 \pm \sqrt{-5}.
            \]
            Now,
            $3 = (\pm 2 \pm \sqrt{-5})(z + w\sqrt{-5})$.
            This implies that either $\frac{3}{2 + \sqrt{-5}} = \frac{3(2 - \sqrt{-5})}{9} = \frac{2}{3} + \dotsc$,
            $\frac{2}{3} \notin \Z$.
            Thus the only factorisations of $3$ that can happen here are factor being a unit $\pm 1$.
            Hence $3$ is irreducible in $\Z[\sqrt{-5}]$.
        \end{enumerate}
    \end{solution}
\end{problem}

\begin{problem}[Question $10$]
    \begin{enumerate}[label = (\alph*)]
        \item Let $r$ be a non-zero element of an integral domain $R$.
        Show that the function $f$ from $R$ to $R$ given by $f(x) = rx$ is injective.

        \item Let $X$ be a finite set.
        Show that any injective function $g : X \to X$ must be surjective.

        \item Conclude that every finite integral domain is a field.
    \end{enumerate}

    \begin{solution}
        \begin{enumerate}[label = (\alph*)]
            \item
            $f$ is injective $f(x) - f(y) \implies x = y$.
            Assume $f(x) = f(y)$,
            i.e.,
            $rx = ry$.
            Then $r(x - y) = 0$.
            Since $r \neq 0$ and $R$ is an integral domain is an integral domain,
            $x - y = 0$
            (can also use cancellation).

            \item
            Let $\Im{g}$ be the image/range of $g$.
            $\Im{g} \subseteq X$,
            hence $|\Im{g}| \leq |X|$.
            Now,
            since $g$ is injective,
            $|\Im{g}| = |X|$,
            so $\mathrm{img} = X$,
            $g$ is surjective.

            \item Let $R$ be a finite integral domain.
            An injective map $R \to R$ must be surjective by the previous part,
            hence bijective.
            By the first part $f(x) = rx$ is bijective,
            so has an inverse.
            By surjectivity,
            $\exists x \in R, rx = 1$.
            Thus $x = r ^ {-1} \in R$.
            Thus $R$ is a field.
        \end{enumerate}
    \end{solution}
\end{problem}

\newpage

\section{Problems Class \texorpdfstring{$3$}{}}

\begin{problem}[Question $11$]
    Let $n \in \Z$.
    For any polynomial
    \[
    f(x) = a_0 + a_1x + \dotsc + a_rx ^ r \in \Z[x]
    \]
    we call
    \[
    \bar{f}(x) = \bar{a}_0 + \bar{a}_1x + \dotsc + \bar{a}_rx ^ r \in (\Z / n)[x]
    \]
    its reduction modulo $n$.
    Prove that the map $\Z[x] \to (\Z / n)[x]$ given by $f(x) \mapsto \bar{f}(x)$ is a homomorphism.


    \begin{solution}
        $f(x)$ as in the question.
        Let $g(x) = b_0 + b_1x + \dotsc + b_sx ^ s$,
        $b_i \in \Z$.
        Additivity:
        \[
        f(x) + g(x) = (a_0 + b_0) + (a_1 + b_1)x + \dotsi.
        \]
        Thus
        \[
        \overline{f(x) + g(x)} = (\overline{a_0 + b_0}) + (\overline{a_1 + b_1})x + \dotsi = (\bar{a}_0 + \bar{b}_0) + (\bar{a}_1 + \bar{b}_1)x + \dotsi = \bar{f}(x) + \bar{g}(x).
        \]
        Multiplicativity:
        Expand the product
        \[
        f(x)g(x) = \sum_{k = 0}^{r + s}\left(\sum_{i = 0}^{k}a_ib_{k - i}\right)x ^ k.
        \]
        Thus,
        \[
        \overline{f(x)g(x)} = \sum_{k = 0}^{r + s}\overline{\left(\sum_{i = 0}^{k}a_ib_{k - i}\right)}x ^ k = \sum_{k = 0}^{r + s}\left(\sum_{i}\bar{a}_i\bar{b}_{k - i}\right)x ^ k = \bar{f}(x)\bar{g}(x).
        \]
        Finally,
        $1 \in \Z[x]$ is mapped to $\bar{1} \in (\Z / n)[x]$.
    \end{solution}
\end{problem}

\begin{problem}[Question $3$]
    Show that the following are irreducible in $\Q[x]$.
    \begin{enumerate}[label = (\roman*)]
        \item $7x ^ 2 - 5x + 2$;
        
        \item $3x ^ 3 + x + 5$;

        \item $x ^ 4 + 35x - 21$;

        \item $3x ^ 4 + 3x ^ 3 + 1$.
    \end{enumerate}

    \begin{solution}\phantom{}
        \begin{enumerate}[label = (\roman*)]
            \item\phantom{}

            \item By lecture notes,
            it is irreducible if and only if it doesn't have any roots in $\Q$.
            Rational root test:
            Assume $b / c \in \Q$ is a root.
            $\gcd(b, c) = 1$,
            $c \mid 3$ and $b \mid 5$.
            Thus,
            the possible roots are $\pm 1$,
            $\pm 5$,
            $\pm \frac{1}{3}$,
            $\pm\frac{5}{3}$.

            Check whether any of these are roots:
            E.g. for $x = 5$,
            $3x ^ 3 + x + 5 > 0$ so $5$ is not a root.
            Show actually that none of them is a root,
            hence the polynomial is irreducible.

            \item
            Always check Eisenstein's.
            Here $p = 7$ works.

            \item Try $p = 2$.
            \[
            \overline{3x ^ 4 + 3x ^ 3 + 1} = x ^ 4 + x ^ 3 + \bar{1} \in (\Z / 2)[x].
            \]
            If we can show that the reduction $x ^ 4 + x ^ 3 +  \bar{1}$ is irreducible in $(\Z / 2)[x]$,
            then the original polynomial in $\Z[x]$ is irreducible.
            Why?
            If the original polynomial is reducible then $3x ^ 4 + 3x ^ 3 + 1 = g(x)h(x)$ and since reduction mod $2$ is multiplicative we get $x ^ 4 + x ^ 3 + \bar{1} = \bar{g}(x)\bar{h}(x)$ is reducible.

            Irreducible in $(\Z / 2)[x]$?
            Roots:
            $0 ^ 4 + 0 ^ 3 + \bar{1} \neq \bar{0}$,
            $\bar{1} + \bar{1} + \bar{1} = \bar{1} \neq \bar{0}$.
            No roots.

            The only irreducible polynomial of degree $2$ in $(\Z / 2)[x]$ is $x ^ 2 + x + \bar{1}$.
            Thus $x ^ 4 + x ^ 3 + \bar{1}$ has no quadratic
            (degree $2$)
            factors,
            as $(x ^ 2 + x + \bar{1}) ^ 2 \neq x ^ 4 + x ^ 3 + \bar{1}$.
            Thus $x ^ 4 + x ^ 3 + \bar{1}$ is irreducible,
            thus the original polynomial is irreducible.
        \end{enumerate}
    \end{solution}
\end{problem}

\begin{problem}[Question $8$]
    Let $p \in \Z$ be a prime number.
    The polynomial
    \[
    \Phi_{p}(x) = x ^ {p - 1} + x ^ {p - 2} + \dotsc + x + 1 = \frac{x ^ p - 1}{x - 1}
    \]
    is called the $p$th cyclotomic polynomial,
    Use Eisenstein's criterion to show that $\Phi_{p}(x)$ is irreducible in $\Q[x]$.

    \begin{solution}
        If $\Phi_p(x + 1)$ is irreducible,
        then it will follow that $\Phi_p(x)$ is irreducible because if $\Phi_p(x) = f(x)g(x)$ with $\deg{f}, \deg{g} \geq 1$,
        then $\Phi_{p}(x + 1) = f(x + 1)g(x + 1)$ is not irreducible.

        Now,
        \[
        \Phi_p(x + 1) = \frac{(x + 1) ^ p - 1}{x} = x ^ {p - 1} + \binom{p}{1}x ^ {p - 2} + \binom{p}{2}x ^ {p - 3} + \dotsi + + \binom{p}{p - 1}.
        \]
        All the non-leading coefficients are divisible by $p$.
        Indeed,
        $\binom{p}{i} = \frac{p!}{i!(p - i)!}$ for $1 \leq i \leq p$,
        which for $i \leq p - 1$ has a factor of $p$ in the numerator but not in the denominator
        (as $i < p$,
        $p - i < p$).
        Thus for $i \leq p - 1$,
        $\binom{p}{i}$ is an integer divisible by $p$.
        Moreover $\binom{p}{p - 1} = p$ is not divisible by $p ^ 2$.
        So by Eisenstein,
        using the prime $p$,
        we conclude that $\Phi_{p}(x + 1)$,
        hence $\Phi_{p}(x)$,
        is irreducible in $\Q[x]$.
    \end{solution}
\end{problem}

\newpage

\section{Problems Class \texorpdfstring{$4$}{}}

\begin{problem}[Question $12$ Problem Sheet $3$]
    Show that the map
    \[
    \varphi : \Z / 6 \to \Z / 2 \times \Z / 3,\qquad\varphi(\bar{a}) = (\bar{a}, \bar{a})
    \]
    is an isomorphism of rings.

    \begin{solution}
        $\Z / 6$ has a subset $A = \{\bar{0}, \bar{2}, \bar{4}\}$ a ring,
        and also another subset $B = \{\bar{0}, \bar{3}\}$ a ring.

        Every element in $\Z / 6$ is uniquely build from element in $A$ and $B$:
        \[
        \bar{0} = \bar{0} + \bar{0},\quad \bar{1} = \bar{4} + \bar{3},\quad \dotsc,\quad \bar{5} = \bar{2} + \bar{3}.
        \]
        $A \cong \Z / 3$,
        $B \cong \Z / 2$.

        Now consider $\varphi$.
        First check $\varphi$ is well-defined,
        i.e.,
        independent of the choices of representatives.
        Assume that $[a]_{6} = [a']_{6}$ for some $a, a' \in \Z$.
        $\varphi([a]_{6}) = ([a]_{2}, [a]_{3})$,
        $\varphi([a']_{6}) = ([a']_{2}, [a']_{3})$
        \begin{align*}
            [a]_{6} = [a']_{6} &\iff a \equiv a'\pmod{6} \\
            &\iff 6 \mid (a - a') \\
            &\iff \begin{cases}
                2 \mid (a - a') \iff [a]_{2} = [a']_{2} \\
                3 \mid (a - a') \iff [a]_{3} = [a']_{3}.
            \end{cases}
        \end{align*}
        Homomorphism:
        $\varphi(\bar{1}) = (\bar{1}, \bar{1})$ identity in $\Z / 2 \times \Z / 3$,
        \[
        \varphi(\bar{a} + \bar{b}) = (\bar{a} + \bar{b}, \bar{a} + \bar{b}) = (\bar{a}, \bar{a}) + (\bar{b} + \bar{b}) = \varphi(\bar{a}) + \varphi(\bar{b}).
        \]
        Injectivity:
        \begin{align*}
            \ker{\varphi} &= \{[a]_{6} \in \Z / 6 \mid ([a]_{2}, [a]_{3}) = (\bar{0}, \bar{0})\} \\
            &= \{[a]_{6} \mid \underbrace{a \equiv 0\pmod{2} \text{ and } a \equiv 0\pmod{3}}_{\iff 2 \mid a\text{ and } 3 \mid a \iff 6 \mid a}\} \\
            &= \{[a]_{6} \mid 6 \mid a\} = \{[\bar{0}]_{6}\}
        \end{align*}
        ($a = 2x = 3y$).

        Surjectivity:
        Note that $|\Z / 6| = 6 = |\Z / 2 \times \Z / 3|$.
        So $\varphi$ is an injective function from a $6$-element set to another $6$-element set.
        Thus $\varphi$ is an isomorphism.
    \end{solution}
\end{problem}

\begin{problem}[Question $13$ Problem Sheet $4$]
    Let $R = \Z[i]$ be the Gaussian integers and consider the ideal $I = (2)$.

    \begin{enumerate}[label = (\roman*)]
        \item Show that $R / I$ has exactly four elements:
        $\bar{0}$,
        $\bar{1}$,
        $\bar{i}$ and $\overline{1 + i}$.

        \item Give the tables for addition and multiplication in $R / I$.

        \item Is $R / I$ isomorphic to $\Z / 2 \times \Z / 2$ or $\Z / 4$ as a ring?
    \end{enumerate}

    \begin{solution}\phantom{}
        \begin{enumerate}[label = (\roman*)]
            \item
            Let $\alpha = a + bi \in \Z[i]$,
            and write $a = 2k + a'$,
            $a' \in \{0, 1\}$,
            some $k \in \Z$.
            $b = 2l + b'$,
            $b' \in \{0, 1\}$,
            some $l \in \Z$.

            Then
            \[
            \bar{\alpha} = \overline{a + bi} = \overline{2k + a'} + \overline{(2l + b')}i = \overline{2k} + \bar{a'} + \overline{2li} + \bar{b'i} = \bar{a'} + \bar{b'i}.
            \]
            Thus there are at most $4$ possible elements in $R / I$.

            Exactly $4$:
            show $\bar{0}, \bar{1}, \bar{i}, \overline{1 + i}$ are distinct.
            \[
            \bar{0} = \bar{1}, \bar{i}, \overline{1 + i}
            \]
            is not true as none of the latter lie in $I$.
            If $\bar{1} = \bar{i}$,
            then $1 - i \in I$ which is not true,
            $\bar{1} = \overline{1 + i}$,
            then $\bar{i} \in I$ which it doesn't.
            If $\bar{i} = \overline{1 + i}$,
            then $\bar{i} \in I$ which is false.

            \item

            \item In $\Z / 4$,
            $\bar{x} + \bar{x}$ is not always $\bar{0}$.
            Thus $R / I$ is not isomorphic to $\Z / 4$.
            This argument is not enough for $\Z / 2 \times \Z / 2$.
            In this direct product $x ^ 2 = x$,
            for all $x \in \Z / 2 \times \Z / 2$ but this doesn't hold for $R / I$.
        \end{enumerate}
    \end{solution}
\end{problem}

\newpage

\section{Problems Class \texorpdfstring{$5$}{}}

\begin{problem}[Sheet $4$ Question $10$]
    Show that if $R$ is a PID and $a \in R$ is irreducible,
    then $(a)$ is a maximal ideal.

    \begin{solution}
        Assume that $I \supseteq (a)$ is an ideal containing $(a)$.
        Want:
        $I = R$ or $I = (a)$..
        Since $R$ is a PID,
        $I = (t)$,
        for some $t \in R$.
        $(t) \supseteq (a) \implies t \mid a$,
        i.e.,
        $a = tm$,
        some $m \in R$.

        $a$ is irreducible so $t$ is a unit or $m$ is a unit.
        If $t$ is a unit,
        then $I = (1) = R$.
        If $m$ is a unit,
        then $(a) = (tm) = (t) = I$.
    \end{solution}
\end{problem}

\begin{problem}[Sheet $4$ Question $11$]
    Show that $\Q[i\sqrt{3}] = \{a + bi\sqrt{3} \mid a, b \in \Q\}$ is a field that is isomorphic to
    \[
    \Q[x] / (x ^ 2 + x + 1).
    \]

    \begin{solution}
        Might be tempted to consider the evaluation map $f(x) \mapsto f(i\sqrt{3})$.
        This does not work,
        as it would not map $x ^ 2 + x + 1$ to $0$,
        hence the kernel would not be $(x ^ 2 + x + 1)$,
        so cannot apply FIT.

        Instead use an evaluation map $\varphi : \Q[x] \to \C$,
        $f(x) \mapsto f(\alpha)$,
        where $\alpha$ is any root of $x ^ 2 + x + 1$
        (e.g.,
        $\alpha = (-1 - i\sqrt{3}) / 2$).

        $\ker{\varphi}$:
        Let $f(x) \in (x ^ 2 + x + 1)$.
        Then $f(x) = (x ^ 2 + x + 1)g(x)$,
        some $g(x) \in \Q[x]$.
        Thus $\varphi(f(x)) = 0$.
        $\varphi(g(x)) = 0$,
        so $(x ^ 2 + x + 1) \subseteq \ker{\varphi}$.

        If $f(x) \in \ker{\varphi}$,
        i.e.,
        $f(\alpha) = 0$.
        Division with remainder
        \[
        f(x) = q(x)(x ^ 2 + x + 1) + r(x),\qquad \deg{r} < 2.
        \]
        Can write $r(x) = ax + b$,
        $a, b \in \Q$,
        \[
        0 = f(\alpha) = 0 + r(\alpha) = a\alpha + b.
        \]
        If $a \neq 0$,
        then $\alpha = -b / a$;
        impossible as $\alpha \notin \Q$.
        Hence $a = 0$.
        Thus $b = 0$.
        This means that $r(x) = 0$,
        so $f(x) \in (x ^ 2 + x + 1)$ so $\ker{\varphi} \subseteq (x ^ 2 + x + 1)$.

        $\Im{\varphi}$:
        Let $f(x) \in \Q[x]$ and write $f(x) = q(x)(x ^ 2 + x + 1) + ax + b$.
        Then $\varphi(f(x)) = a\alpha + b$.
        Using $\alpha = (-1 - i\sqrt{3}) / 2$ this is $\varphi(f(x)) = (b - a / 2) - (a / 2)i\sqrt{3} \in \Q[i\sqrt{3}]$.
        To obtain some $\alpha + \beta i\sqrt{3} \in \Q[i\sqrt{3}]$,
        we only have to solve for $a$ and $b$:
        \[
        b - a / 2 = \alpha
        \]
        \[
        -a / 2 = \beta
        \]
        can be solved,
        so $\Im{\varphi} = \Q[i\sqrt{3}]$.
        Thus by FIT,
        $\Q[x] / (x ^ 2 + x + 1) \cong \Q[i\sqrt{3}]$.
        Since $x ^ 2 + x + 1$ is irreducible in $\Q[x]$
        (no roots in $\Q$),
        Q$10$ says that $\Q[x] / (x ^ 2 + x + 1)$
        (lecture:
        $R / I$ with $I$ maximal,
        is a field).
    \end{solution}
\end{problem}

\begin{problem}[Sheet $4$ Question $14$]
    Show that $(\Z / 5)[x] / I$ where $I = (x ^ 2 + x + \bar{2})$ is a field,
    and find the
    (multiplicative)
    inverse of $\bar{2}x + \bar{3} + I$.

    \begin{solution}
        Field:
        We need $x ^ 2 + x + \bar{2}$ to be irreducible in $(\Z / 5)[x]$.
        Check that no element in $\Z / 5$ is a root.
        Q$10$ implies $(x ^ 2 + x + \bar{2})$ is maximal,
        hence $(\Z / 5)[x] / (x ^ 2 + x + \bar{2})$.

        The inverse:
        By lectures,
        any element in $(\Z / 5)[x] / (x ^ 2 + x + \bar{2})$ is of the form $\bar{a}x + \bar{b} + I$,
        for some $\bar{a}, \bar{b} \in \Z / 5$.
        \begin{align*}
            \bar{a}x + \bar{b} + I= (\bar{2}x + \bar{3} + I) ^ {-1} &\iff (\bar{a}x + \bar{b})(\bar{2}x + \bar{3}) + I = \bar{1} + I \\
            &\iff (\bar{a}x + \bar{b})(\bar{2}x + \bar{3}) - \bar{1} \in I = (x ^ 2 + x + \bar{2}) \\
            &\iff (\bar{a}x + \bar{b})(\bar{2}x + \bar{3}) - \bar{1} = (x ^ 2 + x + \bar{2})g(x),\text{ some $g(x) \in (\Z / 5)[x]$} \\
            &\iff \bar{2}\bar{a}x ^ 2 + (\bar{3}\bar{a} + \bar{2}\bar{b})x + \bar{3}\bar{b} - \bar{1} = (x ^ 2 + x + \bar{2})g(x).
        \end{align*}
        Comparing degrees,
        we must have that $g(x) = c \in \Z / 5$.
        Then,
        \begin{align*}
            \bar{2}\bar{a} &= c \\
            \bar{3}\bar{a} + \bar{2}\bar{b} &= c \\
            \bar{3}\bar{b} - \bar{1} &= \bar{2}c
        \end{align*}
        solving the system,
        we get $\bar{a} = \bar{3}$,
        $\bar{b} = \bar{1}$.
        Thus the inverse of $\bar{2}x + \bar{3} + I$ is $\bar{3}x + \bar{1} + I$.
    \end{solution}
\end{problem}
















\end{document}